\section{A closed source distinguishes two parsing policies}
\label{app:policy}

The distinction between expression values and literal input trees is
visible in one small original source. Let policies L and R agree at
every site except the active three-operand sum partial: L chooses
$(a+b)+1$ and R chooses $a+(b+1)$. Repeats in both are left combs.
Consider
\[
 A=\lambda z:o.\lambda x:o.((\lambda w:o.z)x):o\to o\to o.
\]
This source is closed, recursor-free and has internal type level at
most one. No beta normalization is performed.

Inside assignment trees write $z,x$ for the scalar components $z_0,x_0$.
The assignment of $\lambda w.z$ is $(\ell(z),1)$, with literal tree
\[
 \ell(z)=((z+1)+(z+1))+z,\qquad s_z(\ell)=9.
\]
The body of the $x$ abstraction has scalar
$g=\operatorname{Psi}(1,\ell(z)+x)$, of $x$-size five, so its
abstraction multiplier is ten. Put
\[
 u=\ell+1,\quad T_L=(u+1)+1,\quad T_R=u+(1+1),\quad
 K_\rho=\operatorname{Psi}(1,T_\rho),\quad
 B_* =\operatorname{Psi}(1,u).
\]
The $x$ abstraction has
\[
 h_0=\operatorname{Repeat}_{10}(K_\rho)+B_*,\qquad
 h_1=(1+1)+\operatorname{Sum3}_{\rho}(1,1,1).
\]
The latter is wholly $z$-free and has value five. Direct use of the
size clauses gives
\[
\begin{array}{c|rrrrrrr}
 &\ell&u&T_\rho&K_\rho&\operatorname{Repeat}_{10}(K_\rho)&B_*&h_0\\\hline
 s_z\text{ under L}&9&11&15&17&179&13&193\\
 s_z\text{ under R}&9&11&13&15&159&13&173.
\end{array}
\]
At each scalar tree in this table the first positive output of its
$z$ partial has the corresponding size as its finite value.
Whole-free scalar inputs and a tracked variable both contribute one;
active sums add the child outputs and one; a collapse with head one
adds two to the tail output. This proves the assertion by induction
through the displayed trees, rather than by a numerical evaluation of
the collapse.

The positive output of $\partial_0^zh_0$ is therefore 193 or 173.
The diagonal $(\partial_1^zh_1)_1$ is six, and the retained field
$h_1$ is five. The outer abstraction gives
\[
 \can A_{1,L}=193+6+5=204,\qquad
 \can A_{1,R}=173+6+5=184.
\]
Thus the assigned ordinal values differ, not just their trees.

The scalars differ too. Write $H(t)=\psi_f(\omega+t)$ and
$P_0q=(\partial_0^zq)_0$. The scalar clauses give
$P_0\ell=11$, $P_0u=14$, $P_0T_L=20$ and $P_0T_R=18$.
The two outer size multipliers are 388 and 348; the common retained
scalar is $10H(8)+H(6)$. Consequently
\begin{align*}
 \can A_{0,L}&=388(10H(20)+H(14)+10)+10H(8)+H(6),\\
 \can A_{0,R}&=348(10H(18)+H(14)+10)+10H(8)+H(6).
\end{align*}
Fixed-tail strict increase gives $H(20)>H(18)$, so the L scalar is
strictly larger. This example works for every control in
Section~\ref{sec:setting}. The paper's uniformity in a supplied policy
therefore concerns valid constructions for each policy, not equality
of the resulting assignments across policies.
